\documentclass[11pt]{article}

\usepackage[margin=0.92in]{geometry}
\usepackage{amsmath,amssymb,amsthm,mathtools,booktabs,enumitem,microtype}
\usepackage{xcolor}
\usepackage[colorlinks=true,linkcolor=blue!50!black,urlcolor=blue!50!black,citecolor=blue!50!black]{hyperref}

\theoremstyle{plain}
\newtheorem{theorem}{Theorem}
\newtheorem{lemma}[theorem]{Lemma}
\newtheorem{proposition}[theorem]{Proposition}
\newtheorem{corollary}[theorem]{Corollary}

\theoremstyle{definition}
\newtheorem{definition}[theorem]{Definition}

\theoremstyle{remark}
\newtheorem{remark}{Remark}

\DeclareMathOperator{\Fw}{Fw}
\DeclareMathOperator{\rank}{rank}
\DeclareMathOperator{\spn}{span}

\title{\bfseries Exact Risk Ratios for Weighted Scalar Regression\\on the Upper Half of Budgets}
\author{}
\date{\vspace{-1.6em}%
\small
\begin{tabular}{@{}c@{}}
Target paper: G.~Zhang, \emph{Exact Risk Ratios for Weighted Data Selection in Linear Regression},\\
\href{https://arxiv.org/abs/2608.28007}{arXiv:2608.28007} (October 2026), Conjecture~\texttt{conj:gamma}\\[0.35em]
Companion verification script: \texttt{solutions/weighted\_data\_scalar/explore.py}\\[0.35em]
November 2026
\end{tabular}}

\begin{document}

\maketitle

\begin{abstract}
We prove Conjecture~\texttt{conj:gamma} of Zhang (\href{https://arxiv.org/abs/2608.28007}{arXiv:2608.28007}, October~2026) across the entire upper half of sample budgets $\lfloor d/2 \rfloor \le k \le d - 1$ for every dimension $d \ge 2$:
\[
\Fw(d, d + k) \;=\; 1 + \Gamma_{d, k} \;=\; 2 - \frac{k}{d},
\]
with no support restriction and under the original minimum-Euclidean-norm tie rule.
In particular, $\Fw(5, 7) = 8/5$ and $\Fw(5, 8) = 7/5$.
The proof combines whitened positive-circuit completion and orthogonal-line reduction (settling $2k \ge d$ by a direct induction on dimension) with a structural $3a$-plane covering theorem for rank-$\le 2$ positive-circuit systems that closes the odd boundary $d = 2k + 1$.
The lower half $1 \le k < \lfloor d/2 \rfloor$ remains open; for its first unresolved cell we prove $11/5 \le \Fw(5, 6) \le 12/5$.
\end{abstract}

\section{Setting and Exact Value of \texorpdfstring{$\Gamma_{d,k}$}{Gamma\_\{d,k\}}}

A finite multiset of rows $\{(x_i, y_i)\}_{i=1}^N \subset \mathbb{R}^d \times \mathbb{R}$ has full squared loss $L(w) = N^{-1}\sum_{i=1}^N (x_i^\top w - y_i)^2$ and optimum $L^* = \min_{w \in \mathbb{R}^d} L(w)$.
A weighted selector assigns nonnegative weights to at most $n = d + k$ distinct row indices, and the learner outputs the minimum-Euclidean-norm minimizer of the weighted loss.
The instance risk ratio is the infimum of $L(w)/L^*$ over all such selections (with $0/0 = 1$ and positive$/0 = \infty$), and $\Fw(d, n)$ is the supremum over all finite datasets.

Writing $d = qs + a$ ($0 \le a < s$), Zhang~\cite{Zhang2026} defines
\[
\Gamma_{d, k} := \max_{k < s \le d} \frac{s - k}{(s - a)/q + a/(q + 1)}
\]
and conjectures $\Fw(d, d + k) = 1 + \Gamma_{d, k}$ for all $1 \le k \le d - 1$ (Conjecture~\texttt{conj:gamma}).
Prior unrestricted results~\cite{Hanneke2025,Dewasurendra2026,Zhang2026} covered only $k = d - 1$ and the isolated cells $(3, 4)$, $(4, 5)$, and $(4, 6)$.

\begin{proposition}[Evaluation of $\Gamma_{d,k}$ and lower bound]\label{prop:lb}
For every $d \ge 2$ and $\lceil d/3 \rceil \le k \le d - 1$ (in particular, all $\lfloor d/2 \rfloor \le k \le d - 1$),
\[
\Gamma_{d, k} = \frac{d - k}{d} \qquad\text{and}\qquad \Fw(d, d + k) \ge 2 - \frac{k}{d}.
\]
\end{proposition}

\begin{proof}
For every integer $r \ge 1$, $(r - 1)(r - 2) \ge 0$ gives $1/r \ge (3 - r)/2$.
Summing over any partition of $d$ into $s > k$ positive integer parts yields denominator at least $(3s - d)/2 > 0$.
Because $3k \ge d$ and $s \le d$, $(3k - d)(d - s) \ge 0$ gives $2(s - k)/(3s - d) \le (d - k)/d$, with equality at $s = d$, proving $\Gamma_{d, k} = (d - k)/d$.

For the lower bound, take the $2d$ rows with feature $e_j$ and labels $M \pm 1$ ($j = 1, \dots, d$), where $M := d + 1$.
Then $w^* = M\mathbf{1}$, $L^* = 1$, and excess loss is $d^{-1}\|w - M\mathbf{1}\|_2^2$.
Under the minimum-norm rule, omitting coordinate $j$ costs $M^2/d > 1$, selecting one row costs $1/d$, and selecting both permits zero cost.
With budget $d + k$ and every coordinate covered, at least $d - k$ coordinates receive a single row, yielding ratio $\ge 2 - k/d$.
\end{proof}

\section{Whitened Reductions and Positive Circuits}

Work in $\spn\{x_i\}$, of dimension $m \le d$ (where the minimum-norm rule takes values).
Zero rank gives ratio~$1$; $L^* = 0$ is solved by $m$ independent rows.
Otherwise center at $w^*$ and whiten using $H := \sum x_i x_i^\top \succ 0$: put $z_i := H^{-1/2} x_i$, $t_i := y_i - x_i^\top w^*$, and $u := H^{1/2}(w - w^*)$, so
\[
\sum z_i z_i^\top = I_m, \qquad \sum t_i z_i = 0, \qquad
\sum (z_i^\top u - t_i)^2 = E + \|u\|_2^2 \qquad \Bigl(E := \sum t_i^2\Bigr).
\]
An \emph{excess coefficient} $c$ means an infimum of $\|u\|_2^2 \le c E$ over \emph{full-rank} selections (which have unique minimizers, avoiding tie-rule ambiguity).
Zero-feature rows may be dropped for an upper bound, and we normalize $E = 1$ when $E > 0$.
Write $g_i := t_i z_i$ for the signed gradients.

\begin{lemma}[Flat completion, single rows, and orthogonal lines]\label{lem:tools}
\begin{enumerate}[label=\textup{(\arabic*)}]
\item \emph{Positive circuits and flat completion.} A \emph{positive circuit} is a minimal support of a positive linear dependence among nonzero gradients $g_i$ (its size is its rank plus~$1$; cf.~\cite{Steinitz1916,Davis1954}). A rank-$r$ circuit positively spans its $r$-plane $P$, and weighting its $r + 1$ rows by the dependence coefficients produces a quadratic minimized precisely on $P^\perp$; the same holds for any direct-sum union of circuits spanning $P$. Projecting remaining features onto $P^\perp$ (dropping zero projections) and keeping labels gives a whitened problem on $P^\perp$ with zero gradient sum and energy equal to the sum off $P$. Any full-rank quotient selection lifts to a full-rank selection in $\mathbb{R}^m$ with the same infimal norm bound by multiplying the completed objective by a penalty $\lambda \to \infty$ after choosing finite inner weights.
\item \emph{Anchors and closest hyperplane.} A nonzero row with $t_i = 0$ (an \emph{anchor}) uses $1$ row to fix one coordinate to $0$, reducing $(m, k)$ to $(m - 1, k)$. More generally, since $\sum \|z_i\|_2^2 = m$, the closest interpolation hyperplane $\{u : z_{i_0}^\top u = t_{i_0}\}$ has squared distance $h = t_{i_0}^2/\|z_{i_0}\|_2^2 \le E/m$, attained at $u_0 = (t_{i_0}/\|z_{i_0}\|_2^2)z_{i_0}$. Centering on $z_{i_0}^\perp$ at $u_0$ yields projected gradient sum $\Pi_{z_{i_0}^\perp}\sum (t_i - z_i^\top u_0)z_i = \Pi_{z_{i_0}^\perp}(-u_0) = 0$ and residual energy $\le E + h$. If the $(m - 1)$-dimensional quotient has excess coefficient $D$, lifting gives excess $\le h + D(E + h)$, i.e., ratio factor $\le (1 + 1/m)(1 + D)$. Iterating $m$ times yields ratio $\le m + 1$ with $m$ rows.
\item \emph{Orthogonal lines.} Without anchors, nonzero gradients positively span $\mathbb{R}^m$ and each lies in a positive circuit. If all positive circuits have rank~$1$, each $g_i$ has an opposite partner and the distinct lines are independent (any minimal dependence among distinct lines could be oriented into a higher-rank positive circuit). Writing $I_m = \sum_{j=1}^m Q_j v_j v_j^\top$ (unit $v_j$), the square matrix with columns $\sqrt{Q_j}\,v_j$ is orthogonal, so the $m$ lines are mutually orthogonal with $Q_j = 1$. On line $j$ of residual energy $e_j$, averaging $t_i^2/(z_i^\top v_j)^2$ with weights $(z_i^\top v_j)^2$ gives a single interpolation of cost $\le e_j$, while two opposite rows achieve zero cost. Pairing the $k$ heaviest lines and interpolating the rest yields excess $\le (m - k)E/m$.
\item \emph{Exact recovery for $k \ge m$.} Repeatedly fixing an anchor or completing a positive circuit in the quotient uses at most $2$ rows per removed dimension, achieving zero excess in the infimum with at most $2m \le m + k$ rows.
\end{enumerate}
\end{lemma}

\section{Upper Half \texorpdfstring{$2k \ge d$}{2k >= d}}

\begin{theorem}[Upper half $2k \ge d$]\label{thm:upper}
For every $d \ge 2$ and $\lceil d/2 \rceil \le k \le d - 1$, $\Fw(d, d + k) = 2 - k/d$.
\end{theorem}

\begin{proof}
We establish the full-rank excess coefficient $1 - k/d$ by induction on $d \ge 1$ (base $d = 1$ and $k \ge d$ follow from Lemma~\ref{lem:tools}(4)).
If an anchor exists, Lemma~\ref{lem:tools}(2) reduces to $(d - 1, k)$, improving the coefficient (or yielding exact recovery if $k \ge d - 1$).
If all circuits have rank~$1$, Lemma~\ref{lem:tools}(3) gives $1 - k/d$.
Otherwise complete a circuit of rank $r \ge 2$ (Lemma~\ref{lem:tools}(1)), leaving dimension $m = d - r$ and $k' = k - 1$ extras.
Since $2k \ge d$ and $r \ge 2$, $2k' \ge m$; either $k' \ge m$ (exact recovery) or induction gives
\[
1 - \frac{k - 1}{d - r} \;\le\; 1 - \frac{k}{d} \quad\Longleftrightarrow\quad kr \ge d,
\]
which holds because $r \ge 2$ and $2k \ge d$.
If an input dataset has feature rank $m < d$, the extra budget in $\spn\{x_i\}$ is $k' = d + k - m > k$, giving $k'/m > k/d$.
Combining with Proposition~\ref{prop:lb} completes the proof.
\end{proof}

\section{Rank-\texorpdfstring{$\le 2$}{<= 2} Geometry and the Odd Boundary \texorpdfstring{$d = 2k + 1$}{d = 2k + 1}}

Assume now that there are no anchors and every positive circuit has rank at most~$2$ (call a rank-$2$ positive circuit a \emph{triangle}).
Choose a maximal family of triangles whose planes $P_1, \dots, P_a$ form a direct sum $U := P_1 \oplus \cdots \oplus P_a$, and set $b := d - 2a$.

\begin{lemma}[Quotient lines and $3a$-plane cover]\label{lem:cover}
\begin{enumerate}[label=\textup{(\arabic*)}]
\item Modulo $U$, the nonzero projected gradients lie on $b$ linearly independent lines with both signs.
\item Partition the rows outside $U$ into $b$ groups by their quotient line. Each group is either a single line in $\mathbb{R}^d$ (which is an orthogonal direct summand of the whitened configuration) or is covered by at most $3a$ planes, each carrying an original positive triangle. Moreover, for $a = 1$, a non-single-line group yields a $3$-plane cover of all data in the $3$-space $P_1 \oplus \text{group}$.
\end{enumerate}
\end{lemma}

\begin{proof}
(1) Lifting a positive circuit of rank $s \ge 2$ in $\mathbb{R}^d / U$ with its positive coefficients produces a sum $v \in U$.
If $v = 0$, the lifted circuit has rank $s$ and meets $U$ trivially, contradicting either the rank-$\le 2$ bound ($s \ge 3$) or maximality of $a$ ($s = 2$).
If $v \neq 0$, expressing $-v$ positively by an independent subset of the chosen triangles and appending those columns yields a positive circuit of rank $\ge s + 1 \ge 3$, contradiction.
Hence every quotient circuit has rank~$1$, and Lemma~\ref{lem:tools}(3) yields $b$ independent two-sided quotient lines.

(2) Fix a group and a representative $v \notin U$ of its quotient direction.
Rescaling positively gives gradients $v + p$ ($p \in A \subset U$) and $-(v + q)$ ($q \in B \subset U$) for nonempty finite sets $A, B \subset U$.
For any $(p, q) \in A \times B$, the sum is $p - q \in U$; representing $q - p$ positively by a minimal subset of the chosen triangles uses at most $1$ vector (else a rank-$\ge 3$ positive circuit arises).
Thus every $q - p$ is either $0$ or lies on one of the $3a$ positive rays of the chosen triangles, and each nonzero difference yields a triangle plane containing both rows and that ray.

If $A = B$ is a singleton, the group is a single original line $L$, while all other rows lie in a complementary $(d - 1)$-space $Q$.
The covariance splits as $M_L + M_Q = I_d$ with $\rank(M_L) = 1$ and $\rank(M_Q) = d - 1$, forcing $L \perp Q$, so costs and energies add.
If one of $A, B$ is a singleton and some difference is nonzero, at most $3a$ triangle planes through that row cover the group.

For $a \ge 2$ with $|A|, |B| \ge 2$: if $|A|, |B| \le 2$, at most $4 \le 3a$ unequal pairs cover the group.
Otherwise swap $A$ and $B$ if needed so $|A| \ge 3$; note that negating differences $q - p \mapsto p - q$ preserves the property that every difference lies in a single summand $P_r$, i.e., Hamming distance $d_H(p, q) \le 1$ in $P_1 \times \cdots \times P_a$.
Two points of $B$ differing in $2$ coordinates have only $2$ common neighbors at distance $\le 1$ (and $\ge 3$ coordinates have none), contradicting $|A| \ge 3$.
Hence any two distinct points of $B$ differ in one coordinate $P_r$, forcing all of $A$ (and symmetrically all of $B$) to agree outside $P_r$.
All differences therefore use only the $3$ rays of $P_r$.
Picking an unequal center pair $(p_0, q_0) \in A \times B$, the triangle planes through $p_0$ or $q_0$ number at most $3 + 3 - 1 = 5 \le 3a$ (their shared plane counted once) and cover the group.

For $a = 1$, work in the $3$-space $W := P_1 + \mathbb{R}v$.
A nonzero pair difference yields a second triangle sharing a ray with the first; normalize them to $c + x_1 + x_2 = c + y_1 + y_2 = 0$, which positively span $W$.
Any gradient $g \in W$ outside these two planes has $-g$ positively spanned by at most $2$ independent vectors from $\{c, x_1, x_2, y_1, y_2\}$, so $g = -\alpha x_i - \beta y_j$ ($\alpha, \beta > 0$).
Suppose two cross types occur with $j_1 \neq j_2$ (if only $i_1 \neq i_2$, swap $x$ and $y$).
Rescaling the two vectors $\tilde{g}_1, \tilde{g}_2$ to equalize their $y_{j_1}, y_{j_2}$ coefficients and substituting $y_1 + y_2 = -c = x_1 + x_2$ gives $\tilde{g}_1 + \tilde{g}_2 + \gamma_1 x_1 + \gamma_2 x_2 = 0$ with $\gamma_1, \gamma_2 > 0$.
Since $x_1, x_2$ span the $2$-plane $P_1$ and $\tilde{g}_1 = -\alpha_1 x_{i_1} - \beta_1 y_{j_1}$ ($\beta_1 > 0$) lies outside $P_1$, we have $\dim\spn\{\tilde{g}_1, x_1, x_2\} = 3$.
Thus this $4$-vector strictly positive dependence has a $1$-dimensional kernel, forming a genuine rank-$3$ positive circuit, contradiction.
Hence at most one additional triangle plane exists, giving a $3$-plane cover of all data in $W$.
\end{proof}

\begin{theorem}[Odd boundary $d = 2k + 1$]\label{thm:odd}
For every odd $d = 2k + 1 \ge 3$, $\Fw(d, d + k) = 2 - k/d = 1 + c_d$, where $c_d := (d + 1)/(2d) = 1 - k/d$.
\end{theorem}

\begin{proof}
We induct on odd $d \ge 3$, normalizing $E = 1$.
An anchor reduces to $(d - 1, k)$ (Theorem~\ref{thm:upper}, excess $\le 1/2 < c_d$).
A circuit of rank $\ge 3$ leaves dimension $\le 2(k - 1)$ with $k - 1$ extras (excess $\le 1/2$ by Theorem~\ref{thm:upper}, or exact recovery).
All rank-$1$ circuits are covered by Lemma~\ref{lem:tools}(3).
Thus Lemma~\ref{lem:cover} applies with $a \ge 1$, $b = d - 2a \ge 1$ odd, and $k = a + (b - 1)/2$.

\emph{Base $d = 3$ ($a = 1, b = 1, c_3 = 2/3$).}
If the quotient group is a single original line of energy $e$, completing the triangle costs $\le e$, while pairing the line and interpolating $2$ rows in its orthogonal plane costs $\le 2(1 - e)$ (Lemma~\ref{lem:tools}(2)); their minimum is $\le 2/3$.
Otherwise Lemma~\ref{lem:cover}(2) gives a $3$-plane cover of all data; some triangle plane holds energy $\ge 1/3$, and completing it plus interpolating one quotient row costs $\le 2/3$.

\emph{Inductive step $d \ge 5$.} Set $h := 1/(kd)$.
\begin{enumerate}[label=\textup{(\roman*)},leftmargin=1.8em]
\item \emph{Heavy triangle plane:} If some triangle plane has energy $\ge h$, completing it leaves $(d - 2, k - 1)$ with cost $\le (1 - h)c_{d-2} = c_d$. Henceforth every triangle plane has energy $< h$, so every triangle-covered group has energy $\le 3ah$.
\item \emph{Light single line:} If a single-original-line group has energy $e \le 1/d$, interpolate it and apply the even midpoint (Theorem~\ref{thm:upper}) on its orthogonal complement for cost $\le e + (1 - e)/2 \le c_d$.
\item \emph{Two single lines:} If at least two single-original-line groups remain, their energy sum is $s > 2/d$. Pairing the heavier, interpolating the lighter, and applying odd induction on $(d - 2, k - 1)$ costs $\le s/2 + (1 - s)c_{d-2} \le c_d$ (equality at $s = 2/d$, decreasing in $s$ since $c_{d-2} > 1/2$).
\item \emph{Completing all $a$ triangles:} At most one single-original-line group remains. Complete all $a$ chosen triangles with $3a$ rows, leaving budget $b + (b - 1)/2$ on $b$ orthogonal quotient lines. If $b \ge 3$ (spending one pair on the exceptional line if present) or if $b = 1$ with no single-original-line group, all $(b + 1)/2$ interpolated groups have energy $\le 3ah$. Their total cost is at most
\[
\frac{b + 1}{2}\cdot 3ah \;\le\; c_d
\quad\Longleftrightarrow\quad
4a^2 - 2ab - 6a + b^2 - 1 \;\ge\; 0.
\]
For $b = 1$, this is $4a(a - 2) \ge 0$ (since $d = 2a + 1 \ge 5 \implies a \ge 2$); for $b = 3$, it is $4(a - 1)(a - 2) \ge 0$ for integer $a$; for $b \ge 5$, the quadratic in $a$ has leading coefficient $4 > 0$ and negative discriminant $4(-3b^2 + 6b + 13) < 0$.
\item \emph{Remaining case $b = 1$ with a single original line of energy $e$:} If $e \le c_d$, completing all $a$ triangles and interpolating the line costs $\le e \le c_d$. Otherwise ($e > c_d$), pair the line; on its $(d - 1)$-dimensional orthogonal complement apply one closest-hyperplane step (Lemma~\ref{lem:tools}(2)) followed by odd induction on $(d - 2, k - 1)$. The complement excess coefficient is at most
\[
D := \left(1 + \frac{1}{d - 1}\right)(1 + c_{d-2}) - 1 = \frac{d^2 + d - 4}{2(d - 2)(d - 1)} \;\le\; \frac{d + 1}{d - 1}
\]
for $d \ge 5$ (cross-multiplied difference $d(d - 3) > 0$), giving total cost $\le D(1 - e) \le D(1 - c_d) \le c_d$.
\end{enumerate}
Rank-deficient inputs ($m < d$) have $k' = d + k - m \ge \lfloor m/2 \rfloor$ and $1 - k'/m \le 1 - k/d$.
Combining with Proposition~\ref{prop:lb} finishes the proof.
\end{proof}

\section{First Unresolved Cell \texorpdfstring{$(5, 6)$}{(5, 6)} and Verification}

\begin{proposition}[Bracket for $(5, 6)$]\label{prop:56}
$11/5 \le \Fw(5, 6) \le 12/5$.
\end{proposition}

\begin{proof}
For the upper bound on full-rank data, apply one closest-hyperplane reduction (Lemma~\ref{lem:tools}(2), factor $6/5$) followed by Zhang's Theorem~\texttt{thm:fw45}~\cite[Appendix~F]{Zhang2026} (which furnishes full-rank / strictly convex selections for $(4, 5)$ with ratio $\le 2$), yielding $\Fw(5, 6) \le (6/5) \cdot 2 = 12/5$; rank-deficient inputs ($m \le 4$) have budget $6 \ge m + 2$ and satisfy a strictly tighter bound.
For the lower bound, take two orthogonal simplex blocks of ranks $r = 2$ and $r = 3$, each with $r + 1$ rows $e_1, \dots, e_r, -\sum_{j=1}^r e_j$ and residuals $a_r > 0$ satisfying $r(r + 1)a_r^2 = C$, shifted by $Mv$ with $v$ outside all proper row spans (forcing rank-deficient selections to diverge as $M \to \infty$).
With $6$ rows at most one block is completed; breaking the other costs $C$ against optimal loss $C(1/2 + 1/3) = \frac{5}{6}C$, yielding ratio $1 + 6/5 = 11/5$.
\end{proof}

\begin{remark}[Scope caveat]
Theorems~\ref{thm:upper} and~\ref{thm:odd} settle Conjecture~\texttt{conj:gamma} for all $d \ge 2$ and $\lfloor d/2 \rfloor \le k \le d - 1$.
The lower half $1 \le k < \lfloor d/2 \rfloor$ (starting at $(5, 6)$, where Proposition~\ref{prop:56} leaves a gap between $11/5$ and $12/5$) remains open.
\end{remark}

The companion script \texttt{solutions/weighted\_data\_scalar/explore.py} verifies $\Gamma_{d, k} = (d - k)/d$ across $3{,}333$ cells ($d \le 100$), checks all $5{,}049$ odd-packing inequalities ($d \le 201$), exhaustively tests the shifted-axis extremizer through $d = 9$, runs $33$ coupled graph-data experiments across dimensions $3$--$6$, and tests $96$ coupled $2 + 3$ compositions for $(5, 6)$ (largest observed excess $0.91619392$, versus $1.2$ for two simplices). All assertions pass.

\begin{thebibliography}{9}

\bibitem{Davis1954}
C.~Davis.
Theory of positive linear dependence.
\emph{American Journal of Mathematics}, 76(4):733--746, 1954.

\bibitem{Dewasurendra2026}
M.~Dewasurendra.
\emph{Risk Ratios for Weighted Linear Regression}.
Zenodo, 2026. \href{https://doi.org/10.5281/zenodo.21881535}{doi:10.5281/zenodo.21881535}.

\bibitem{Hanneke2025}
S.~Hanneke, S.~Moran, A.~Shlimovich, and A.~Yehudayoff.
Weighted data selection for linear regression.
In \emph{Proceedings of the 38th Annual Conference on Learning Theory (COLT)}, PMLR~291, 2025.

\bibitem{Steinitz1916}
E.~Steinitz.
Bedingt konvergente Reihen und konvexe Systeme.
\emph{Journal f\"ur die reine und angewandte Mathematik}, 146:1--52, 1916.

\bibitem{Zhang2026}
G.~Zhang.
\emph{Exact Risk Ratios for Weighted Data Selection in Linear Regression}.
Preprint \href{https://arxiv.org/abs/2608.28007}{arXiv:2608.28007}, October 2026.

\end{thebibliography}

\end{document}
